\section{引言}

在此开始填入完整中文正文。保持与原文一致的章节层级、公式编号、脚注和交叉引用。
译文应以顺畅、清晰的学术中文为准，可按中文表达习惯调整语序，必要时拆分长句或合并短句，但不得增删论证内容。
专业术语首次出现时请写作“中文术语（English Term）”；之后保持译名或缩写一致，不必在每次出现时重复括注。

插图示例：
% \begin{figure}[htbp]
%   \centering
%   \includegraphics[width=0.92\linewidth]{fig-1.png}
%   \caption{图 1：翻译后的图注。图内文字保留原文。}
%   \label{fig:fig1}
% \end{figure}

表格图片示例：
% \begin{figure}[htbp]
%   \centering
%   \includegraphics[width=\linewidth]{table-1.png}
%   \caption{表 1：翻译后的表注。表内文字保留原文图片。}
%   \label{fig:table1}
% \end{figure}
